\documentclass[conference]{IEEEtran}
\usepackage{cite}
\usepackage{graphicx}
\usepackage{booktabs}
\usepackage{amsmath}
\usepackage[hidelinks]{hyperref}

% Final report: IEEE double-column, up to 8 pages (course requirement).
% Builds on proposal/main.tex. references.bib starts as a copy of proposal/references.bib.

\begin{document}

\title{PhishLang-Web: Auditing and Deploying Fully In-Browser Phishing Detection with MobileBERT}

\author{\IEEEauthorblockN{Masroor Ul Hasan Syed \qquad Deepthi Kondreddy}
\IEEEauthorblockA{CS 8027: Advanced Networking \& Architecture, Fall 2026}}

\maketitle

\begin{abstract}
% TODO: write last, from the results.
\end{abstract}

\section{Introduction}
% Problem, why client-side, contributions list (RQ1--RQ3 answers).

\section{Background and Related Work}
% PhishLang~\cite{roy2026phishlang}, MobileBERT~\cite{sun2020mobilebert}, LLM detectors, client-side anti-phishing, evaluation pitfalls.

\section{Threat Model and Problem Statement}

\section{Audit of PhishLang's Evaluation (RQ1--RQ2)}
% Data, leakage-free and temporal splits, context-window ablation.

\section{In-Browser Deployment (RQ3)}
% ONNX export, INT8 quantization, JS parser, MV3 architecture.

\section{Evaluation}
% Accuracy parity, latency, memory, package size across machines.

\section{Discussion and Limitations}

\section{Conclusion and Future Work}

\nocite{*} % TODO: remove once the text cites its references (it lists every bib entry)
\bibliographystyle{IEEEtran}
\bibliography{references}

\end{document}
